\thispagestyle{plain}
\addcontentsline{toc}{chapter}{Abstract}

\begin{center}
	{\Large \bf ABSTRACT}
\end{center}

Generative AI is exploding in popularity, and it’s really changed how we think about whether digital media is real. The image makers we have now, using diffusion, Generative Adversarial Networks or GANs, and systems that work with multiple kinds of data to create images, can make photos that look absolutely real to people. This makes it easy for false information, scams, and a decline in faith in anything we see online to spread around. As for checking if an image is real, the tools we have at the moment - older Convolutional Neural Networks, and the more demanding Vision-Language Models - become much less accurate with pictures that have a lot going on, are in a particular artistic style, or have words on them. The industry has largely compensated for this deficiency through multi-model ensemble architectures, however, these systems are computationally prohibitive, resource-intensive, and fundamentally too slow for real-time web deployment. This dissertation presents a novel, lightweight, single-pass hybrid AI image detector that matches the robustness of heavy ensembles while operating at a fraction of the latency. The proposed architecture isolates the spatial feature extraction capabilities of the Sigmoid Loss for Language-Image Pre-Training (SigLIP) 2 vision encoder by deliberately removing its text-processing components, thereby reducing memory footprint and inference cost. To capture fine-grained forensic details invisible at standard resolutions, the model's position embeddings are bicubically interpolated from $224 \times 224$ to $512 \times 512$. These rich deep embeddings are fused with seven explicitly computed traditional frequency and statistical forensic signals---Local Binary Patterns (LBP), Discrete Cosine Transform (DCT), Gradient Co-occurrence, Haar Wavelet Energy, Fast Fourier Transform (FFT) Power Slope, Bayer Noise Pattern, and Edge Sharpness Grid---which collectively address the blind spots of deep learning alone. A custom three-layer classification head employing Squeeze-and-Excitation (SE) style channel attention is then applied to the fused embedding to identify AI-generated artefact patterns. Evaluated on a stratified 15\% held-out test split of the training dataset (1,132 images), the model achieves 96.6\% accuracy, 98.9\% precision, 95.3\% recall, 97.0\% F1-score, and a 0.992 AUC-ROC, with an inference latency of 142.4 milliseconds measured on a dual NVIDIA Tesla T4 GPU environment on Kaggle Notebook. A comparative analysis against commercial detection platforms Undetectable AI (5,230 ms) and AI or Not (920 ms) demonstrates that our architecture correctly identifies hard-negative AI-generated images misclassified as real by both commercial tools, while achieving a $6.5\times$ and $36.7\times$ speed advantage respectively. The model is deployed as an interactive Hugging Face Spaces application, providing accessible, real-time image verification.

\begin{keywords}
	AI Image Detection, Synthetic Media, Vision-Language Models, SigLIP 2, Feature Fusion, Image Authentication.
\end{keywords}

\clearpage
